\thispagestyle{empty}

\begin{center}
	\begin{tabular}{lcl}
		\multirow{5}{*}{\includegraphics[height=3cm]{portada/ipn.jpg}} &
		&
      	\multirow{3}{*}{\includegraphics[height=2.8cm]{portada/escom.jpg}}\\
      	& \large \bf INSTITUTO POLITÉCNICO NACIONAL & \\
      	& \large \bf ESCUELA SUPERIOR DE CÓMPUTO &\\
      	& \large \bf SUBDIRECCIÓN ACADÉMICA &
	\end{tabular}
\end{center}

\vspace{0.5cm}

\centerline{\large No. de Registro: 2014-A021}
\centerline{\large Documento técnico}

\vspace{0.8cm}

\centerline{\Large \bf Traveler Assistant System For Mexico City (TASMC)}

\vspace{0.8cm}

\centerline{\large Presentan}

\vspace{0.3cm}

\centerline{\large \bf Barajas Uribe Sergio \footnote{E-mail: scscf.1992@gmail.com}}
\centerline{\large \bf Vivanco Carmona Erick Rafael \footnote{E-mail: erickvivanco01@hotmail.com}}

\vspace{0.5cm}

\centerline{\large Directores}

\vspace{0.3cm}

\centerline{\large \bf M. en C. Macario Hernández Cruz}
\centerline{\large \bf M. en C. Axel Ernesto Moreno Cervantes}

\vspace{0.8cm}

\centerline{\Large Resumen}

El presente proyecto TASMC, consta del desarrollo de un sistema de información que tendrá como propósito proporcionar al usuario elementos necesarios para la organización integral del viaje aéreo, con carácter turístico o de negocios, en el Aeropuerto Internacional de la Ciudad de México (AICM), suministrando en un dispositivo móvil funciones que serán desarrolladas en dos fases: la primera corresponde a la ubicación del usuario dentro del aeropuerto; la segunda fase corresponde a la construcción de bases de información para seleccionar la mejor opción  de vuelos, hoteles, sugerencia de rutas al aeropuerto a través de mapas, información de vuelos en lo que respecta a salidas y llegadas.

\vspace{0.2cm}

\centerline{\textbf{Palabras Clave }- Aplicación Móvil, Geolocalización, Aplicación Web e Ingeniería de Software}